\chapter{Automorphic Forms and the Decomposition of \texorpdfstring{$\bL^2(\Gamma\backslash \SL(2,\R))$}{L2(Gamma SL(2,R))}}

\label{ch:2}

Let $\bG$ denote the group $\SL(2,\R)$ and $\Gamma$ the congruence subgroup $\Gamma_0(N)$. One purpose of this section is to describe the correspondence between $S_k(\Gamma)$ and a certain finite-dimensional subspace of $\bL^2(\Gamma \setminus \bG)$. This correspondence motivates a quite general notion of ``automorphic form'' which includes as special cases the non-analytic ``wave-forms'' of Maass as well as the classical holomorphic forms just discussed.

Another purpose of this section is to explain the role automorphic forms play in decomposing the natural representation of $\bG$ in $\bL^2(\Gamma \setminus \bG)$ and to collect some diverse facts concerning this decomposition.

\section{Automorphic Forms as Functions on \texorpdfstring{$\SL(2,\R)$}{SL(2,R)}}\label{sec:2a}

In $\bG$, consider the subgroups

\[
\mathbf{A} = \left\{\begin{pmatrix} a & 0 \\ 0 & a^{-1} \end{pmatrix}: a > 0\right\}, \quad \mathbf{N} = \left\{\begin{pmatrix} 1 & u \\ 0 & 1 \end{pmatrix}\right\},
\]

and

\[
\bK = \left\{\begin{pmatrix} \cos\theta & -\sin\theta \\ \sin\theta & \cos\theta \end{pmatrix} = r(\theta): 0 \leq \theta < 2\pi\right\}.
\]

The subgroup $B = \mathbf{N}\mathbf{A}$ already acts transitively on $\{\operatorname{Im}(z) > 0\}$ since $\begin{pmatrix} y^{1/2} & xy^{-1/2} \\ 0 & y^{-1/2} \end{pmatrix} i = x + iy$. Thus the upper half-plane is identified with $\bG/\bK$, the stability subgroup of $\bG$ at $i$ being $\bK$.

Since $\bG = \mathbf{B}\bK = \mathbf{N}\mathbf{A}\bK$, each $g \in \bG$ may be expressed in the form

\begin{equation}\label{eq:2.1}
g = \begin{pmatrix} a & b \\ c & d \end{pmatrix} = \begin{pmatrix} y^{1/2} & xy^{-1/2} \\ 0 & y^{-1/2} \end{pmatrix} \begin{pmatrix} \cos\theta & -\sin\theta \\ \sin\theta & \cos\theta \end{pmatrix}.\tag{2.1}
\end{equation}

Hence, assigning to each $g$ of this form the coordinates $(z = x + iy, \theta)$, a convenient parameterization of $\bG$ (by $x, y, \theta$), $x \in \R$, $y > 0$, $\theta \in [0, 2\pi])$ is obtained. Observe that

\begin{equation}\label{eq:2.2}
z = g(i), \quad \text{and} \quad \theta = \arg(ci + d),\tag{2.2}
\end{equation}

if $g = \begin{pmatrix} a & b \\ c & d \end{pmatrix}$.

Now we describe a map from $S_k(\Gamma)$ to a space of functions on $\bG$. For simplicity we assume $k$ is even. Then for each $f \in S_k(\Gamma)$ we define $\phi_f(g)$ on $\bG$ by

\begin{equation}\label{eq:2.3}
\phi_f(g) = f(g(i)) j(g,i)^{-k}.\tag{2.3}
\end{equation}

Observe that (by the transitivity of $\bG$) the function $\phi_f(g)$ is identically zero if and only if $f(z)$ is. The resulting map, $f$ goes to $\phi_f$, is a linear one-to-one map from $S_k(\Gamma)$ to a space of functions $\phi$ on $\bG$ satisfying the following conditions:

(i) $\phi(\gamma g) = \phi(g)$ for all $\gamma \in \Gamma$;

(ii) $\phi(gr(\theta)) = e^{-ik\theta}\phi(g)$ for all $r(\theta) \in \bK$;

(iii) $\phi(g)$ is bounded, in particular

\begin{equation}\label{eq:2.4}
\int_{\Gamma \setminus \bG} |\phi(g)|^2 \dd g < \infty;\tag{2.4}
\end{equation}

and

(iv) $\phi(g)$ is cuspidal, i.e. for $g \in \bG$ and $\sigma \in \SL(2,\Z)$

\[
\int_0^1 \phi\left(\sigma \begin{pmatrix} 1 & xh \\ 0 & 1 \end{pmatrix} g\right) \dd x = 0.
\]

(Here $h$ denotes the ``width'' of the cusp $\sigma(\infty)$ in the sense of Section~\ref{ch:1}.)

Condition 2.4(i) follows from the computation

\[
\begin{aligned}
\phi_f(\gamma g) &= f(\gamma g(i)) j(\gamma g, i)^{-k}\\
&= j(\gamma, g(i))^k f(g(i)) j(\gamma, g(i))^{-k} j(g, i)^{-k} = \phi_f(g).
\end{aligned}
\]

(Recall that $j$ is a factor of automorphy, i.e. $j(g_1g_2, z) = j(g_1, g_2z) j(g_2, z)$ for all $g_1, g_2 \in \bG$.) Similarly (ii) follows from the computation

\[
\begin{aligned}
\phi_f(gr(\theta)) &= f(gr(\theta)(i)) j(gr(\theta), i)^{-k}\\
&= f(g(i)) j(g, i)^{-k} j(r(\theta), i)^{-k} = \phi_f(g) e^{-ik\theta}.
\end{aligned}
\]

(Observe that the restriction of $j(g,i)^{-k}$ to $\bK$ is the character $e^{-ik\theta}$.)

On the other hand, (iii) follows immediately from Lemma~\ref{lem:1.5}, and

\[
\begin{aligned}
\int_0^1 \phi_f\left(\sigma \begin{pmatrix} 1 & xh \\ 0 & 1 \end{pmatrix} g\right) \dd x &= \int_0^1 f\left(\sigma \begin{pmatrix} 1 & xh \\ 0 & 1 \end{pmatrix} g(i)\right) j\left(\sigma \begin{pmatrix} 1 & xh \\ 0 & 1 \end{pmatrix} g, i\right)^{-k} \dd x\\
&= \int_0^1 f(\sigma(z + hx)) j(\sigma, z + hx)^{-k} \dd x\\
&= \int_0^1 f\big|_{[\sigma]_k}(hx + z) \dd x.
\end{aligned}
\]

Since this last expression is simply the zeroth Fourier coefficient of $f$ at the cusp $s = \sigma(\infty)$ (cf. \eqref{eq:1.8}) Condition 2.4(iv) follows from the fact $f$ is a cusp form in the classical sense.

Perhaps a word is in order concerning the Haar measure used on $\Gamma \setminus \bG$ in 2.4(iii). Utilizing the parameterization $(x, y, \theta)$ in $\bG$ we normalize Haar measure on $\bG$ through the formula

\begin{equation}\label{eq:2.5}
\int_{\bG} \phi(g) \dd g = \frac{1}{2\pi} \int_0^{2\pi} \int_0^{\infty} \int_{-\infty}^{\infty} \phi(x, y, \theta) \frac{\dd x \dd y}{y^2} \dd\theta.\tag{2.5}
\end{equation}

Consequently

\begin{equation}\label{eq:2.6}
\int_{\Gamma \setminus \bG} |\phi_f(g)|^2 \dd g = \iint_F |f(z)|^2 y^k \frac{\dd x \dd y}{y^2}\tag{2.6}
\end{equation}

where $F$ denotes a fundamental domain for $\Gamma$ in the upper half-plane. The right side of \eqref{eq:2.6} is simply the norm of $f(z)$ with respect to the Petersson inner product.

The question remains: what is the image of $S_k(\Gamma)$ in $\bL^2(\Gamma \setminus \bG)$ under the map $f \to \phi_f$? To answer this we need to recall the Laplace operator for $\bG$.

Let $\bR(g)$ denote the unitary representation of $\bG$ defined in $\bL^2(\Gamma \setminus \bG)$ by the formula

\begin{equation}\label{eq:2.7}
\bR(g)\phi(h) = \phi(hg).\tag{2.7}
\end{equation}

(This is the so-called right regular representation of $\bG$ in $\bL^2(\Gamma \setminus \bG)$ induced from the trivial representation of $\Gamma$.) Consider the matrices

\[
\ell_0 = \begin{pmatrix} 0 & -1 \\ 1 & 0 \end{pmatrix}, \quad \ell_1 = \begin{pmatrix} 0 & 1 \\ 1 & 0 \end{pmatrix},
\]

and

\[
\ell_2 = \begin{pmatrix} 1 & 0 \\ 0 & -1 \end{pmatrix}
\]

which span the Lie algebra of $\SL(2,\R)$. To each one-parameter group $g_j(t) = \exp(t\ell_j)$ in $\bG$ there corresponds the one-parameter group of unitary operators $U_j(t) = \bR(g_j(t))$, hence, by Stone's Theorem, a self-adjoint (unbounded) operator $H_j$ in $\bL^2(\Gamma \setminus \bG)$ satisfying $U_j(t) = e^{-itH_j}$.

The Laplace operator for $\bG$ in $\bL^2(\Gamma \setminus \bG)$ (i.e. the Casimir operator for $\bR(g)$) is

\begin{equation}\label{eq:2.8}
\Delta = -\frac{1}{4}\left(H_0^2 - H_1^2 - H_2^2\right).\tag{2.8}
\end{equation}

This operator is defined on smooth functions in $\bL^2(\Gamma \setminus \bG)$ and in terms of the coordinates $(x, y, \theta)$ is easily seen to assume the form

\begin{equation}\label{eq:2.9}
\Delta = -y^2\left(\frac{\partial^2}{\partial x^2} + \frac{\partial^2}{\partial y^2}\right) - y \frac{\partial^2}{\partial x \partial \theta};\tag{2.9}
\end{equation}

Moreover, (the minimal closed extension of) $\Delta$ is self-adjoint in $\bL^2(\Gamma \setminus \bG)$ and commutes with $\bR(g)$. Consequently it follows from Schur's Lemma that the restriction of $\Delta$ to any $\bG$-invariant irreducible subspace of $\bL^2(\Gamma \setminus \bG)$ is a scalar.

We can now characterize the image of $S_k(\Gamma)$ in $\bL^2(\Gamma \setminus \bG)$. Let $A_k^2(\Gamma)$ denote the space of functions $\phi$ on $\bG$ satisfying the following conditions:

(a) $\phi(\gamma g) = \phi(g)$ for all $\gamma \in \Gamma$;

(b) $\phi(gr(\theta)) = \phi(g)e^{-ik\theta}$ for all $r(\theta) \in \bK$;

(2.10) (c) $\Delta\phi = -\frac{k}{2}\left(\frac{k}{2} - 1\right)\phi$; and

(d) $\phi$ is bounded and cuspidal.

Note that $A_k^2(\Gamma)$ is contained in $\bL^2(\Gamma \setminus \bG)$ by 2.10(d).

\begin{proposition}\label{prop:2.1}

The formula

\[
\phi_f(g) = f(g(i)) j(g, i)^{-k}
\]

describes an isomorphism between $S_k(\Gamma)$ and $A_k^2(\Gamma)$.

\end{proposition}

\begin{proof}

We have only to show that $\phi_f$ satisfies 2.10(c) and that the map is onto (cf. 2.4). So first suppose $f \in S_k(\Gamma)$. Since $\phi_f(g) = y^{k/2}f(z)e^{-ik\theta} = \phi_f(z, \theta)$, a straightforward computation shows that

\begin{equation}\label{eq:2.11}
\begin{aligned}
\Delta\phi_f(g) &= \left(-y^2\left(\frac{\partial^2}{\partial x^2} + \frac{\partial^2}{\partial y^2}\right) - y \frac{\partial^2}{\partial x \partial \theta}\right)\phi(z, \theta)\\
&= -e^{-ik\theta} y^{(k/2) + 2}\left(\frac{\partial^2}{\partial x^2} + \frac{\partial^2}{\partial y^2}\right) f(z)\\
&\quad + (ik)e^{-ik\theta} y^{k/2 + 1}\left(\frac{\partial}{\partial x} + i\frac{\partial}{\partial y}\right) f(z)\\
&\quad - \frac{k}{2}\left(\frac{k}{2} - 1\right) y^{k/2} f(z) e^{-ik\theta}.
\end{aligned}\tag{2.11}
\end{equation}

But $f(z)$ is holomorphic by assumption. Therefore the first two terms on the right side of \eqref{eq:2.11} vanish and 2.10(c) is established.

Now suppose $\phi(g)$ is arbitrary in $A_k^2(\Gamma)$ and put

\begin{equation}\label{eq:2.12}
f(z) = \phi(g) j(g, i)^{k}.\tag{2.12}
\end{equation}

(Here $g$ is any element at all of $\bG$ such that $g(i) = z$.) From 2.10(b) it follows that $f(z)$ is well-defined. On the other hand $\phi_f$ certainly equals $\phi$. Therefore to prove this proposition it remains only to check that $f_{\phi}(z) \in S_k(\Gamma)$. Since this verification requires facts not yet introduced we shall simply sketch the argument for the sake of completeness.

Let $\bL_0^2(\Gamma \setminus \bG)$ denote the space of $\phi$ in $\bL^2(\Gamma \setminus \bG)$ satisfying 2.4(iv) for almost every $g \in \bG$. By hypothesis, $\phi \in \bL_0^2(\Gamma \setminus \bG)$. On the other hand (as we shall see in Theorem~\ref{thm:2.6}(ii)),

\[
\bL_0^2(\Gamma \setminus \bG) = \bigoplus_j \bH^{j}
\]

where each $\bH^{j}$ is invariant for the right action of $\bG$ and irreducible. Consequently

\[
\phi = \sum \phi_{\ell}
\]

where $\phi_{\ell} \in \bH^{\ell}$.

Note now that each $\phi_j$ satisfies 2.10(b) since $\phi$ does. Recalling that $\Delta$ operates as a scalar on each irreducible subspace of $\bL^2(\Gamma \setminus \bG)$, we also have

\[
\Delta\phi_{\ell} = -\frac{k}{2}\left(\frac{k}{2} - 1\right)\phi_{\ell}
\]

for each $\ell$, since

\[
\lambda_{\ell}(\phi, \phi_{\ell}) = (\phi, \Delta\phi_{\ell}) = (\Delta\phi, \phi_{\ell}) = -\frac{k}{2}\left(\frac{k}{2} - 1\right)(\phi, \phi_{\ell})
\]

(by hypothesis and the self-adjointness of $\Delta$). From this it may be deduced that $\phi$ is a linear combination of ``lowest weight vectors'' from irreducible subspaces of $\bL_0^2(\Gamma \setminus \bG)$ equivalent to the discrete series representation $\pi_k^{+}$ (see Section B). Moreover, each such vector (and hence $\phi$ itself) can be shown to be annihilated by a certain first order differential operator in $(x, y, \theta)$ which in turn implies that $f_{\phi}(z)$ is annihilated by the first order operator

\[
\frac{\partial}{\partial x} + i\frac{\partial}{\partial y}.
\]

Consequently $f_{\phi}(z)$ is holomorphic. (The details of this argument appear in Section 4 of Chapter 1 of \cite{GelfandGraevPyatetskiiShapiro}.)

That $f_{\phi}(z)$ is automorphic of weight $k$ and a cusp form follows immediately by reversing the computations used earlier to show that $\phi_f \in A_k^2(\Gamma)$ whenever $f \in S_k(\Gamma)$.

\end{proof}

According to Proposition~\ref{prop:2.1}, classical (holomorphic) cusp forms are examples of right K-finite eigenfunctions of $\Delta$ in $\bL^2(\Gamma \setminus \bG)$. (Right K-finite means that the space of functions on $\bG$ spanned by the right translates of $\phi(g)$ by $k$ in $\bK$ is finite-dimensional.) Now from the point of view of the spectral decomposition of $\Delta$ there is nothing particularly special about the eigenvalue $-k/2(k/2 - 1)$ (cf. 2.10(c)). Therefore it is natural to make the following definition.

\begin{definition}\label{def:2.2}

A $\Gamma$-automorphic form $\phi$ on $\bG$ is a (smooth) function on $\bG$ satisfying the following properties:

(a) $\phi(\gamma g) = \phi(g)$ for all $\gamma \in \Gamma$;

(b) $\phi$ is right K-finite;

(c) $\phi$ is an eigenfunction of $\Delta$; and

(d) $\phi$ satisfies a certain growth condition:

there are constants $C$ and $N$ such that

\begin{equation}\label{eq:2.13}
|\phi(z, \theta)| \leq Cy^{N} \quad (y \to +\infty).\tag{2.13}
\end{equation}

(We say that $\phi$ is slowly increasing.)

If, in addition to satisfying properties 2.13(a)-(d), $\phi$ is cuspidal in the sense of 2.4(iv), then $\phi$ is called a $\Gamma$-cusp form. Of course elements of $S_k(\Gamma)$ (or rather $A_k^2(\Gamma)$) constitute the fundamental examples of $\Gamma$-automorphic forms on $\bG$. Others include:

\end{definition}

\begin{example}[\normalfont (Non-holomorphic cusp forms.)]\label{exa:2.3}

Let $W_s(\Gamma)$ denote the space of smooth right K-invariant functions $\phi$ on $\bG$ which are bounded and such that

\[
\Delta\phi = \frac{1 - s^2}{4}\phi.
\]

Then each $\phi$ in $W_s(\Gamma)$ is obviously an automorphic form on $\bG$ in the sense of Definition~\ref{def:2.2}. Furthermore, the formula

\[
f_{\phi}(z) = \phi(g(i)),
\]

where $g$ in $\bG$ maps $i$ to $z$, establishes an isomorphism between $W_s(\Gamma)$ and the finite-dimensional space of complex-valued functions $f(z)$ on $\{\operatorname{Im}(z) > 0\}$ satisfying the following properties:

(i) $f(\gamma z) = f(z)$ for all $\gamma \in \Gamma$ (i.e. $f(z)$ is an automorphic function!);

(2.14) (ii) $\Delta^{*}f = \frac{1 - s^2}{4}f$, where $\Delta^{*} = -y^2\left(\frac{\partial^2}{\partial x^2} + \frac{\partial^2}{\partial y^2}\right)$ is the Laplace-Beltrami operator for the upper half-plane ($f$ is real analytic, not holomorphic);

(iii) $f(z)$ is bounded.

These functions were first systematically studied by H. Maass and called by him ``wave forms.''

\end{example}

\begin{remark}\label{rem:2.4}

In Condition (ii) of Example~\ref{exa:2.3} we must assume $s$ to be a pure imaginary number or a non-zero real number between $-1$ and $1$. This is because the Laplace-Beltrami operator is non-negative.

\end{remark}

Soon we shall see that $W_s(\Gamma)$ is non-trivial for infinitely many values of $s$. However, to the best of our knowledge, almost nothing is known concerning those specific values of $s$ for which $W_s(\Gamma)$ is non-trivial. (For certain proper subgroups of $\SL(2,\Z)$, a few examples are known; cf. Theorem~\ref{thm:2.14} below.)

\begin{example}[\normalfont (Real Analytic Eisenstein Series.)]\label{exa:2.5}

Suppose $\Gamma = \SL(2,\Z)$, $\mu \in \C$, and $z = x + iy$, $y > 0$. If $f_{\mu}(z) = y^{\mu}$ the series

\begin{equation}\label{eq:2.15}
E^{*}(z,\mu) = \sum_{\gamma \in \Gamma \cap N \backslash \Gamma} f_{\mu}(\gamma z) = \sum_{\substack{(c,d) = 1 \\ c,d \in \Z}} \frac{y^{\mu}}{|cz + d|^{2\mu}}\tag{2.15}
\end{equation}

is known to converge for $\operatorname{Re}(\mu) > 1$ and to have a meromorphic continuation to the whole $\mu$-plane with a simple pole at $\mu = 1$. The resulting function of $z$ is known as the Eisenstein series of index $\mu$ attached to the cusp $\infty$ of $\Gamma$ and its significance is as follows. If $\mu = 1/2 + s/2$, and $s$ is pure imaginary, then the formula

\begin{equation}\label{eq:2.16}
E(g, s) = E^{*}\left(g(i), \frac{1}{2} + \frac{s}{2}\right)\tag{2.16}
\end{equation}

describes a right K-invariant $\Gamma$-automorphic form on $\bG$ which is not a cusp form. Its eigenvalue for $\Delta$ is $\frac{1 - s^2}{4}$.

\end{example}

Generalizations of $E(g,s)$ are discussed in Section~\ref{ch:8} and play a fundamental role in the theory of automorphic forms. In fact Examples \ref{exa:2.3} and \ref{exa:2.5}, together with $A_k^2(\Gamma)$, essentially exhaust all possible automorphic forms on $\bG$.

\section{Automorphic Forms and the Decomposition of \texorpdfstring{$\bL^2(\Gamma\backslash G)$}{L2(Gamma G)}}\label{sec:2b}

As before, let $\bR(g)$ denote the right regular representation of $\bG$ in $\bL^2(\Gamma \setminus \bG)$.

From Proposition~\ref{prop:2.1} we know that there is a one-to-one correspondence between the elements of $S_k(\Gamma)$ and certain eigenfunctions of $\Delta$ in $\bL^2(\Gamma \setminus \bG)$. Similarly, by Example~\ref{exa:2.3}, there is a one-to-one correspondence between certain other eigenfunctions of $\Delta$ and real-analytic wave-forms.

On the other hand, the restriction of $\Delta$ to irreducible subrepresentations of $\bR(g)$ is known a priori to operate as a scalar. Therefore one must expect automorphic forms to play a role in the decomposition of $\bR(g)$. That they do provides a specific instance of the general principle whereby eigenspaces of an operator commuting with a given representation reduce that representation. In the context of $\bR(g)$ this operator is clearly $\Delta$. And since the restriction of $\bR(g)$ to $\bK$ is completely reducible ($\bK$ being compact) there is certainly no loss of generality in requiring these eigenfunctions to be right K-finite. Therefore the problem of decomposing $\bR(g)$ is virtually indistinguishable from the problem of constructing automorphic forms.

In this subsection we shall directly consider the representation theoretic problem of decomposing $\bR(g)$.

We start by describing the irreducible unitary representations of $\SL(2,\R)$.

(a) The class 1 principal series $\pi_s^{+}(g)$.

These representations are induced from the unitary characters

\begin{equation}\label{eq:2.17}
\begin{pmatrix} a & u \\ 0 & a^{-1} \end{pmatrix} \mapsto |a|^{s} \quad (\operatorname{Re}(s) = 0)\tag{2.17}
\end{equation}

of $B$. Hence they are described as right translation operators in the Hilbert space $\bH^{+}(s)$ consisting of measurable functions $\phi$ on $\bG$ such that

\begin{equation}\label{eq:2.18}
\phi\left(\begin{pmatrix} a & u \\ 0 & a^{-1} \end{pmatrix} g\right) = |a|^{s+1}\phi(g)\tag{2.18}
\end{equation}

and

\begin{equation}\label{eq:2.19}
\int_{\bK} |\phi(k)|^2 \dd k = \|\phi\|^2 < \infty.\tag{2.19}
\end{equation}

We recall that $\pi_s^{+}$ and $\pi_{-s}^{+}$ are unitarily equivalent.

(b) The non-class 1 principal series $\pi_s^{-}(g)$.

These representations are induced from the characters

\[
\begin{pmatrix} a & u \\ 0 & a^{-1} \end{pmatrix} \mapsto \sgn(a)|a|^{s}, \quad \operatorname{Re}(s) = 0,
\]

and are irreducible if and only if $s \neq 0$. Hence we assume $s \neq 0$. Again, $\pi_s^{-}$ and $\pi_{-s}^{-}$ are equivalent.

(c) The complementary series $\pi_s^{c}(g)$.

These unitary representations result from induction from the non-unitary characters

\[
\begin{pmatrix} a & u \\ 0 & a^{-1} \end{pmatrix} \to |a|^{s}
\]

of $B$ with $s$ a non-zero real number between $-1$ and $1$. Of course the inner product in $\bH(s)$ invariant for $\pi_s^{c}(g)$ is longer given by \eqref{eq:2.19} but by a more complicated formula which we omit. Suffice it to say that such an invariant inner product does exist and results from the positive-definiteness of the analytic continuation of the operator intertwining the equivalent representations $\pi_s^{+}$ and $\pi_{-s}^{+}$.

(d) The discrete series $\pi_k^{\pm}(g)$.

These representations are subrepresentations of certain non-unitary representations again induced from $B$ but may be conveniently described as follows. For each integer $k > 1$, let $\bH(k)$ denote the Hilbert space of holomorphic functions in the upper half-plane satisfying

\begin{equation}\label{eq:2.20}
\iint_{\operatorname{Im}(z) > 0} |f(x + iy)|^2 y^k \frac{\dd x \dd y}{y^2} < \infty.\tag{2.20}
\end{equation}

Then

\begin{equation}\label{eq:2.21}
\pi_k^{+}(g) f(z) = (bz + d)^{-k} f\left(\frac{az + c}{bz + d}\right)\tag{2.21}
\end{equation}

if $g = \begin{pmatrix} a & b \\ c & d \end{pmatrix}$. The representations $\pi_k^{-}(g)$ are defined similarly by replacing the upper half-plane by the lower and $y^k$ by $|y|^k$.

(e) The trivial representation.

It is a well-known fact due to Bargmann that the representations just described exhaust the totality of irreducible unitary representations of $\bG$. Therefore each irreducible subrepresentation of $\bR(g)$ is equivalent to some one of these and the question remains which irreducible representations so appear.

Observe that since $\Gamma \setminus \bG$ is non-compact there is no reason to expect the decomposition of $\bR(g)$ to be entirely discrete. Moreover it should not seem surprising that this discrete spectrum is (essentially) exhausted by the cuspidal functions in $\bL^2(\Gamma \setminus \bG)$.

To make matters precise, suppose $\Gamma$ is $\SL(2,\Z)$ and let $\bL_0^2(\Gamma \setminus \bG)$ denote the subspace of $\phi$ in $\bL^2(\Gamma \setminus \bG)$ satisfying the cuspidal condition 2.4(iv). Note that for $\SL(2,\Z)$ this condition is

\[
\int_{\mathbf{N} \cap \Gamma \backslash \mathbf{N}} \phi(ng) \dd n = 0 \quad \text{for a.e.}\ g
\]

Since this condition is obviously invariant under the right action of $\bG$, the space of cusp forms $\bL_0^2(\Gamma \setminus \bG)$ is invariant for $\bR(g)$.

\begin{theorem}\label{thm:2.6}

Let $\bL_d^2(\Gamma \setminus \bG)$ denote the direct sum of all subspaces of $\bL^2(\Gamma \setminus \bG)$ irreducibly invariant under the action of $\bG$. Then:

\begin{enumerate}[(i)]

\item $\bL_d^2(\Gamma \setminus \bG)$ is the direct sum of $\bL_0^2(\Gamma \setminus \bG)$ and the one-dimensional subspace of constant functions;

\item The spectrum of $\bL_0^2(\Gamma \setminus \bG)$ has finite multiplicity, i.e. in the restriction of $\bR(g)$ to $\bL_0^2(\Gamma \setminus \bG)$ each irreducible representation of $\bG$ occurs at most finitely many times;

\item The restriction of $\bR(g)$ to the orthocomplement of $\bL_d^2(\Gamma \setminus \bG)$ is the continuous sum of the principal series representations $\pi_s^{+}(g)$, i.e. if $\bR_c(g)$ denotes the restriction of $\bR(g)$ to the orthocomplement of $\bL_d^2(\Gamma \setminus \bG)$ then

\begin{equation}\label{eq:2.22}
\bR_c(g) = \bigoplus \int_0^{\infty} \pi_s^{+}(g) \dd s.\tag{2.22}
\end{equation}

\end{enumerate}

\end{theorem}

This theorem appears as a corollary to the discussions of Sections~\ref{ch:8} and 9, and, as already remarked, is equivalent to the spectral decomposition of $\Delta$. It's proof involves showing that any slowly increasing function $f(z)$ on $\{\operatorname{Im}(z) > 0\}$ which is orthogonal to all cusp forms may be expressed as an integral average

\[
\int_{\substack{\operatorname{Re}(s) = 0 \\ \operatorname{Im}(s) > 0}} \hat{f}\left(\frac{1}{2} + \frac{s}{2}\right) E(z, s) \dd s
\]

of real analytic Eisenstein series $E(z,s)$.

\begin{remark}\label{rem:2.7}

Suppose $\Gamma'$ is Hecke's subgroup $\Gamma_0(N)$ and $m$ denotes the number of inequivalent cusps of $\Gamma$. Then the continuous sum \eqref{eq:2.22} becomes an $m$-dimensional continuous sum and Parts (i) and (ii) of Theorem~\ref{thm:2.6} remain unchanged. Now suppose $\Gamma'$ is an arbitrary Fuchsian group of the first kind. Then Part (i) is no longer valid since in this generality infinite dimensional representations of $\bG$ may occur in $\bL_d^2(\Gamma \setminus \bG)$ outside the space of cusp forms. Indeed square-integrable $\Gamma$-automorphic forms which are not cusp forms arise as residues at the poles of Eisenstein series (cf. Section~\ref{ch:8}) and for general $\Gamma$ these poles may be more interesting than the pole at 1 is for $\Gamma = \Gamma_0(N)$.

\end{remark}

\begin{remark}\label{rem:2.8}

Although Theorem~\ref{thm:2.6} completely describes the continuous spectrum of $\bR(g)$ it tells us nothing about the discrete spectrum except that this spectrum has finite multiplicities. The question remains which representations of $\bG$ occur discretely and with what multiplicity? The only general result along these lines is a ``duality theorem'' essentially due to Gelfand, Fomin, and Pyatetskii-Shapiro. To describe it we need first to collect several basic facts from the representation theory of $\bG$. (For more complete treatments of these well-known facts the reader is referred to the ``Notes and References'' at the end of this section.)

\end{remark}

\begin{lemma}\label{lem:2.9}

If $\bH^{\pm}(s)$ (resp. $\bH^{c}(s)$) denotes an irreducible subspace of $\bL^2(\Gamma \setminus \bG)$ equivalent to $\pi_s^{\pm}(g)$ (resp. $\pi_s^{c}(g)$) then

\begin{equation}\label{eq:2.23}
\Delta\phi = \frac{1 - s^2}{4}\phi\tag{2.23}
\end{equation}

on $\bH^{\pm}(s)$ (resp. $\bH^{c}(s)$). Similarly

\begin{equation}\label{eq:2.24}
\Delta\phi = -\frac{k}{2}\left(\frac{k}{2} - 1\right)\phi\tag{2.24}
\end{equation}

on $\bH^{\pm}(k)$.

\end{lemma}

\begin{proof}[Proof (Sketch)]

This lemma is equivalent to the statement that the Casimir operator for the irreducible representation $\pi_s^{\pm}$, $\pi_s^{c}$, or $\pi_k^{\pm}$ (cf. \eqref{eq:2.8}) is described by \eqref{eq:2.22} or \eqref{eq:2.23}. But this last fact follows from a straightforward computation using the formulas

\[
H_j = \lim_{t \to 0} \frac{\pi(\exp t\ell_j) - I}{t} = \frac{\dd}{\dd t}\left(\pi(\exp t\ell_j)\right)\big|_{t = 0}
\]

\[
\Delta = -1/4\left(H_0^2 - H_1^2 - H_2^2\right)
\]

and convenient realizations of $\pi_s^{\pm}$ and $\pi_s^{c}$.

\end{proof}

To describe the result of Gelfand, Fomin, Pyatetskii-Shapiro we need to recall the following basic facts from the representation theory of $\bG$:

(i) The restriction to $\bK$ of an irreducible unitary representation $\pi$ of $\bG$ contains a given irreducible representation of $\bK$ at most once;

(ii) This restriction contains the trivial representation of $\bK$ if and only if $\pi$ is some $\pi_s^{+}$ or $\pi_s^{c}$: These representations have precisely one $\bK$-fixed vector in their space and comprise the so-called class 1 representations of $\bG$;

(iii) The restriction of $\pi_k^{+}$ to $\bK$ contains the $\bK$-representation $r(\theta) \to e^{\mp ik\theta}$.

We can now state and prove:

\begin{theorem}\label{thm:2.10}

The representation $\pi_k^{+}$ of $\bG$ (resp. $\pi_s^{+}$ or $\pi_s^{c}$) occurs in $\bR_0(g)$ (the restriction of $\bR(g)$ to $\bL_0^2(\Gamma \setminus \bG)$) with multiplicity equal to the dimension of $S_k(\Gamma)$ (resp. $W_s(\Gamma)$).

\end{theorem}

\begin{proof}

By Theorem~\ref{thm:2.6} we can write

\begin{equation}\label{eq:2.25}
\bR_0(g) = \bigoplus_j \pi^{j}\tag{2.25}
\end{equation}

where $\pi^{j}$ is irreducible and acts on some subspace $\bH^{j}$ of $\bL^2(\Gamma \setminus \bG)$ by right translation. So consider the representation of $\bK$ defined in $\bH^{j}$ by restriction to $\bK$, i.e. by the operators $\pi^{j}(r(\theta))$, $r(\theta) \in \bK$. Since $\bK = \SO(2)$ is compact and abelian,

\[
\bH^{j} = \bigoplus \{\phi_{\ell}^{j}(g)\},
\]

where each $\phi_{\ell}^{j}$ is in the domain of $\Delta$ and such that $\phi_{\ell}^{j}(gr(\theta)) = e^{-i\ell\theta}\phi_{\ell}^{j}(g)$. Now suppose $\bH^{j}$ is equivalent to $\pi_k^{+}$, say. Then exactly one $\phi_{\ell}^{j}$ in $\bH^{j}$, namely $\phi_k^{j}$ (the ``lowest weight vector for $\pi_k^{+}$''), simultaneously satisfies the conditions

\begin{equation}\label{eq:2.26}
\phi(gr(\theta)) = e^{-ik\theta}\phi(g)\tag{2.26}
\end{equation}

and

\begin{equation}\label{eq:2.27}
\Delta\phi = -\frac{k}{2}\left(\frac{k}{2} - 1\right)\phi.\tag{2.27}
\end{equation}

Consequently, by Proposition~\ref{prop:2.1}, exactly one $\phi_{\ell}^{j}$ in $\bH^{j}$ ``belongs'' to $S_k(\Gamma)$. Hence $m_k^{+}$, the multiplicity of $\pi_k^{+}$ in $\bR_0(g)$, is at most the dimension of $S_k(\Gamma)$.

To obtain the reverse inequality we let $f$ denote any non-zero element of $S_k(\Gamma)$. Then $\phi_f(g) = f(g(i))j(g,i)^{-k}$ belongs to $\bL_0^2(\Gamma \setminus \bG)$ and satisfies \eqref{eq:2.26} and \eqref{eq:2.27}, again by Proposition~\ref{prop:2.1}.

By the complete reducibility of $\bR_0(g)$ (cf. 2.24),

\[
\phi_f = \sum_j \phi_j,
\]

with each $\phi_j$ a smooth function in some irreducible subspace $\bH^{j}$ of $\bL_0^2(\Gamma \setminus \bG)$. These $\phi_j$ again satisfy \eqref{eq:2.26}, and $\Delta\phi_j = \lambda_j\phi_j$ for some $\lambda_j$ by the irreducibility of $\bH^{j}$. Actually, by the self-adjointness of $\Delta$, $(\phi, \Delta\phi_j) = (\Delta\phi, \phi_j)$, so $\lambda_j$ must equal $-k/2(k/2 - 1)$. Consequently $\bH^{j}$ must be equivalent to $\pi_k^{+}$. Since this implies that the dimension of $S_k(\Gamma)$ is at most $m_k^{+}$, the proof of the theorem is complete for $\pi_k^{+}$.

The proof for $\pi_s^{\pm}$ and $\pi_s^{c}$ is similar except that in these cases one deals with right K-invariant $\phi_j$ and appeals to Example~\ref{exa:2.3}.

\end{proof}

Observe that if $\begin{pmatrix} -1 & 0 \\ 0 & -1 \end{pmatrix} \in \Gamma$ then $\bR_0(g)$ can never contain the representations $\pi_s^{c}$ or $\pi_k^{+}$ with $k$ odd.

\section{Some Miscellaneous Results Concerning the Decomposition of \allowbreak\texorpdfstring{$\bL^2(\Gamma\backslash \SL(2,\R))$}{L2(Gamma SL(2,R))}}\label{sec:2c}

As far as principal and complementary series are concerned the duality theorem of the last subsection yields very little explicit information. This is because for any fixed value of $s$ very little is known concerning the dimension of $W_s(\Gamma)$.

The following theorem shows that we at least know $W_s(\Gamma)$ is non-empty for infinitely many (mysterious) values of $s$.

\begin{theorem}\label{thm:2.11}

In the decomposition of $\bR_0(g)$, infinitely many class 1 representations occur. Equivalently, $W_s(\Gamma)$ is non-empty for infinitely many permissible values of $s$.

\end{theorem}

\begin{proof}[Proof (Sketch)]

Letting $\bR_0(g)$ denote right translation in the space of cusp forms $\bL_0^2(\Gamma \setminus \bG)$ we again write

\begin{equation}\label{eq:2.28}
\bR_0(g) = \bigoplus_j \pi^{j}\tag{2.28}
\end{equation}

where each $\pi^{j}$ is an irreducible unitary representation of $\bG$. Clearly the subspace of $\bK$-fixed vectors in $\bL_0^2(\Gamma \setminus \bG)$ coincides with the Hilbert space $\bL_0^2(\Gamma \setminus \bG / \bK) = \bL_0^2(\Gamma \setminus \{\operatorname{Im}(z) > 0\})$. However, it can be shown that this latter space is infinite-dimensional. Therefore, since each $\pi^{j}$ in \eqref{eq:2.28} contains the trivial representation of $\bK$ at most once, infinitely many of these $\pi^{j}$ must be class 1.

\end{proof}

This theorem, of course, does not tell us whether representations of both the principal and complementary series occur. More generally, it certainly does not provide specific values of $s$ for which $W_s(\Gamma)$ is non-empty.

\begin{conjecture}\label{conj:2.12}

If $\pi_s^{+}$ belongs to the principal series of representations, then $\pi_s^{+}$ occurs in $\bL^2(\SL(2,\Z) \setminus \bG)$ at most once.

\end{conjecture}

\begin{conjecture}\label{conj:2.13}

No representation at all of the complementary series for $\bG$ occurs in $\bL^2(\SL(2,\Z) \setminus \bG)$.

\end{conjecture}

Interesting experimental evidence supporting the truth of these conjectures has recently been obtained by Cartier. The second conjecture is actually a theorem and is attributed to Selberg (see \cite{Roelcke}).

The following special result is a sample of possible consequences of Jacquet-Langland's theory and its proof is postponed until Section~\ref{ch:7}.

\begin{theorem}\label{thm:2.14}

The principal series representation $\pi_s^{+}$, with

\[
s = \frac{2\pi i}{\log(\sqrt{2} - 1)},
\]

occurs discretely in $\bL^2(\Gamma_0(2) \setminus \bG)$.

\end{theorem}

A last miscellaneous result concerns the discrete spectrum of an arbitrary Fuchsian group of the first kind and was communicated to us by R. Hotta.

Recall that in the generality of an arbitrary Fuchsian group non-trivial representations of $\bG$ can occur discretely outside the space of cusp forms. The thrust of the result below is that these representations never belong to the discrete series for $\bG$. This fact appears naturally as a special consequence of Harish-Chandra's general theory of discrete series and cusp forms for arbitrary semi-simple groups. Thus we state it without proof rather than digress from the subject matter of this section.

\begin{theorem}\label{thm:2.15}

Suppose $\Gamma$ is an arbitrary Fuchsian group of the first kind and $\pi_k$ is a representation of the discrete series with $k > 2$. Then if $\pi_k$ occurs in $\bL_d^2(\Gamma \setminus \bG)$ it occurs in $\bL_0^2(\Gamma \setminus \bG)$.

\end{theorem}

\textbf{Concluding Remark.} Theorem~\ref{thm:2.15} also follows from the theory of Eisenstein series. Indeed this theory implies that each ``non-cuspidal'' subrepresentation of $\bL_d^2(\Gamma \setminus \bG)$ ``imbeds'' in some $\pi_s^{c}(g)$ with $s \in [0,1]$. (Cf. Sections~\ref{ch:4}.A and~\ref{ch:8}.C.)

\section*{Notes and References}

\addcontentsline{toc}{subsection}{Notes and References}

We recommend \cite{GelfandGraevPyatetskiiShapiro} for further discussion of the basic notions of this section (at least for the case of compact quotient). The connection between automorphic forms and group representations seems first to have been made explicit by \cite{GelfandFomin} in 1952. Definition~\ref{def:2.2} is a special case of the general notion of an automorphic form on an arbitrary semi-simple group introduced in \cite{HarishChandra},

The theory of real analytic Eisenstein series was introduced in \cite{Selberg} for Fuchsian groups and subsequently significantly generalized in \cite{Langlands}, A detailed treatment of Selberg's theory (in particular, Theorem~\ref{thm:2.6}) has recently been given in \cite{Kubota}. Further discussion of Eisenstein series and the continuous spectrum is to be found in Section~\ref{ch:8}.

The experimental evidence supporting Conjectures \ref{conj:2.12} and \ref{conj:2.13} appears in \cite{Cartier}.
